% GENERATED FILE. Edit canonical lessons, then run npm run generate:books.
% canonical-source-hash: fnv1a64:16e5f5a1948c072b
% canonical-lessons: HI-C325-angrez, HI-C325-andolan, HI-C325-sanghars, HI-C325-balidan, HI-C325-saumpna

\chapter{A Movement for Freedom}
\label{ch:hi-a-movement-for-freedom}
\hlchaptermodality{325}

\begin{chapteropening}
\textbf{By the end of this chapter:} \emph{I can tell how a freedom movement grew.}

Complete the last lesson of chapter 325: I can tell how a freedom movement grew.
\end{chapteropening}

\section[\texorpdfstring{\hi{अंग्रेज़} (aṅgrez)}{अंग्रेज़ (aṅgrez)}]{\hi{अंग्रेज़} (aṅgrez) — an English person}

\label{lesson:HI-C325-angrez}



\begin{quote}
Before the new one: say the Hindi for a rebellion, then the Hindi for to challenge.
\end{quote}

\subsection*{You'll want to know: \hi{अंग्रेज़}}
\textbf{\hi{अंग्रेज़}} — \emph{aṅgrez} — \textquotedblleft{}an English person\textquotedblright{}.

\begin{quote}
\hi{अंग्रेज़} \hi{दो} \hi{सौ} \hi{साल} \hi{तक} \hi{इस} \hi{देश} \hi{में} \hi{रहे।}
\end{quote}

\begin{quote}
\emph{aṅgrez do sau sāl tak is deś meṁ rahe.}
\end{quote}

\begin{quote}
The English stayed in this country for two hundred years.
\end{quote}

Masculine or feminine, by the person meant.

\begin{cousinweb}
\textbf{\hi{अंग्रेज़}} (\emph{aṅgrez}) is usually traced to Portuguese \emph{inglês}, \textquotedblleft{}English\textquotedblright{}, through Persian \emph{angrez}. \textbf{\hi{अंग्रेज़ी}}, the English language, is built on it.
\end{cousinweb}

\begin{grammarlens}[title={What you've built}]
One more word for telling what happened.
\end{grammarlens}

\subsection*{Your turn}
\begin{itemize}
\raggedright
  \item \emph{Say it:} \emph{aṅgrez}
  \item \emph{Say it:} \emph{aṅgrez}, once more
  \item \emph{Say it:} \emph{aṅgrez do sau sāl tak is deś meṁ rahe.} — The English stayed in this country for two hundred years.
  \item \emph{Recall:} say the Hindi for a rebellion, then the Hindi for to challenge, then say \emph{aṅgrez} again
\end{itemize}

\begin{tcolorbox}[breakable,colback=teal!4,colframe=teal!35!black,title={Before you move on}]
What does \hi{अंग्रेज़} mean? (\textquotedblleft{}An English person\textquotedblright{}.) Say it once more.
\end{tcolorbox}
\section[\texorpdfstring{\hi{आंदोलन} (āndolan)}{आंदोलन (āndolan)}]{\hi{आंदोलन} (āndolan) — a movement (of people)}

\label{lesson:HI-C325-andolan}



\begin{quote}
Before the new one: say the Hindi for to challenge, then the Hindi for an English person.
\end{quote}

\subsection*{You'll want to know: \hi{आंदोलन}}
\textbf{\hi{आंदोलन}} — \emph{āndolan} — \textquotedblleft{}a movement (of people)\textquotedblright{}.

\begin{quote}
\hi{आज़ादी} \hi{का} \hi{आंदोलन} \hi{पूरे} \hi{देश} \hi{में} \hi{फैल} \hi{गया।}
\end{quote}

\begin{quote}
\emph{āzādī kā āndolan pūre deś meṁ phail gayā.}
\end{quote}

\begin{quote}
The freedom movement spread across the whole country.
\end{quote}

Masculine.

\begin{cousinweb}
\textbf{\hi{आंदोलन}} (\emph{āndolan}) is Sanskrit \emph{āndolana}, \textquotedblleft{}a swinging, a rocking\textquotedblright{} (Monier-Williams). Modern Hindi uses it for a political movement; the dictionary has only the older sense.
\end{cousinweb}

\begin{grammarlens}[title={What you've built}]
One more word for telling what happened.
\end{grammarlens}

\subsection*{Your turn}
\begin{itemize}
\raggedright
  \item \emph{Say it:} \emph{āndolan}
  \item \emph{Say it:} \emph{āndolan}, once more
  \item \emph{Say it:} \emph{āzādī kā āndolan pūre deś meṁ phail gayā.} — The freedom movement spread across the whole country.
  \item \emph{Recall:} say the Hindi for to challenge, then the Hindi for an English person, then say \emph{āndolan} again
\end{itemize}

\begin{tcolorbox}[breakable,colback=teal!4,colframe=teal!35!black,title={Before you move on}]
What does \hi{आंदोलन} mean? (\textquotedblleft{}A movement (of people)\textquotedblright{}.) Say it once more.
\end{tcolorbox}
\section[\texorpdfstring{\hi{संघर्ष} (saṅgharṣ)}{संघर्ष (saṅgharṣ)}]{\hi{संघर्ष} (saṅgharṣ) — a struggle}

\label{lesson:HI-C325-sanghars}



\begin{quote}
Before the new one: say the Hindi for an English person, then the Hindi for a movement (of people).
\end{quote}

\subsection*{You'll want to know: \hi{संघर्ष}}
\textbf{\hi{संघर्ष}} — \emph{saṅgharṣ} — \textquotedblleft{}a struggle\textquotedblright{}.

\begin{quote}
\hi{आज़ादी} \hi{के} \hi{लिए} \hi{लोगों} \hi{ने} \hi{लंबा} \hi{संघर्ष} \hi{किया।}
\end{quote}

\begin{quote}
\emph{āzādī ke liye logoṁ ne lambā saṅgharṣ kiyā.}
\end{quote}

\begin{quote}
People waged a long struggle for freedom.
\end{quote}

Masculine.

\begin{cousinweb}
\textbf{\hi{संघर्ष}} (\emph{saṅgharṣ}) is Sanskrit \emph{saṁgharṣa} (Monier-Williams), a rubbing together and so a rivalry: in Hindi, a long fight against hardship or against a ruler.
\end{cousinweb}

\begin{grammarlens}[title={What you've built}]
One more word for telling what happened.
\end{grammarlens}

\subsection*{Your turn}
\begin{itemize}
\raggedright
  \item \emph{Say it:} \emph{saṅgharṣ}
  \item \emph{Say it:} \emph{saṅgharṣ}, once more
  \item \emph{Say it:} \emph{āzādī ke liye logoṁ ne lambā saṅgharṣ kiyā.} — People waged a long struggle for freedom.
  \item \emph{Recall:} say the Hindi for an English person, then the Hindi for a movement (of people), then say \emph{saṅgharṣ} again
\end{itemize}

\begin{tcolorbox}[breakable,colback=teal!4,colframe=teal!35!black,title={Before you move on}]
What does \hi{संघर्ष} mean? (\textquotedblleft{}A struggle\textquotedblright{}.) Say it once more.
\end{tcolorbox}
\section[\texorpdfstring{\hi{बलिदान} (balidān)}{बलिदान (balidān)}]{\hi{बलिदान} (balidān) — a sacrifice}

\label{lesson:HI-C325-balidan}



\begin{quote}
Before the new one: say the Hindi for a movement (of people), then the Hindi for a struggle.
\end{quote}

\subsection*{You'll want to know: \hi{बलिदान}}
\textbf{\hi{बलिदान}} — \emph{balidān} — \textquotedblleft{}a sacrifice\textquotedblright{}.

\begin{quote}
\hi{उस} \hi{आंदोलन} \hi{में} \hi{हज़ारों} \hi{लोगों} \hi{ने} \hi{बलिदान} \hi{दिया।}
\end{quote}

\begin{quote}
\emph{us āndolan meṁ hazāroṁ logoṁ ne balidān diyā.}
\end{quote}

\begin{quote}
Thousands of people gave their lives in that movement.
\end{quote}

Masculine.

\begin{cousinweb}
\textbf{\hi{बलिदान}} (\emph{balidān}) is Sanskrit \emph{balidāna}, \emph{bali}, an offering, and \emph{dāna}, giving (Monier-Williams). In history it is giving one's life for others.
\end{cousinweb}

\begin{grammarlens}[title={What you've built}]
One more word for telling what happened.
\end{grammarlens}

\subsection*{Your turn}
\begin{itemize}
\raggedright
  \item \emph{Say it:} \emph{balidān}
  \item \emph{Say it:} \emph{balidān}, once more
  \item \emph{Say it:} \emph{us āndolan meṁ hazāroṁ logoṁ ne balidān diyā.} — Thousands of people gave their lives in that movement.
  \item \emph{Recall:} say the Hindi for a movement (of people), then the Hindi for a struggle, then say \emph{balidān} again
\end{itemize}

\begin{tcolorbox}[breakable,colback=teal!4,colframe=teal!35!black,title={Before you move on}]
What does \hi{बलिदान} mean? (\textquotedblleft{}A sacrifice\textquotedblright{}.) Say it once more.
\end{tcolorbox}
\section[\texorpdfstring{\hi{सौंपना} (saumpnā)}{सौंपना (saumpnā)}]{\hi{सौंपना} (saumpnā) — to entrust; to hand over}

\label{lesson:HI-C325-saumpna}



\begin{quote}
Before the new one: say the Hindi for a struggle, then the Hindi for a sacrifice.
\end{quote}

\subsection*{You'll want to know: \hi{सौंपना}}
\textbf{\hi{सौंपना}} — \emph{saumpnā} — \textquotedblleft{}to entrust; to hand over\textquotedblright{}.

\begin{quote}
\hi{जाने} \hi{से} \hi{पहले} \hi{राजा} \hi{ने} \hi{क़िले} \hi{की} \hi{चाबी} \hi{अपने} \hi{बेटे} \hi{को} \hi{सौंपी।}
\end{quote}

\begin{quote}
\emph{jāne se pahle rājā ne qile kī cābī apne beṭe ko saumpī.}
\end{quote}

\begin{quote}
Before leaving, the king handed the key of the fort to his son.
\end{quote}

Past, with \hi{ने} (he or she …): \hi{उसने} \hi{चाबी} \hi{सौंपी}, \emph{usne cābī saumpī}.

\begin{cousinweb}
\textbf{\hi{सौंपना}} (\emph{saumpnā}) goes back to Sanskrit \emph{samarpayati}, \textquotedblleft{}hands over\textquotedblright{} (CDIAL 13192). What is handed over agrees with the verb: \textbf{\hi{चाबी} \hi{सौंपी}}; the person who receives it takes \textbf{\hi{को}}.
\end{cousinweb}

\begin{grammarlens}[title={What you've built}]
One more word for telling what happened.
\end{grammarlens}

\subsection*{Your turn}
\begin{itemize}
\raggedright
  \item \emph{Say it:} \emph{saumpnā}
  \item \emph{Say it:} \emph{saumpnā}, once more
  \item \emph{Say it:} \emph{jāne se pahle rājā ne qile kī cābī apne beṭe ko saumpī.} — Before leaving, the king handed the key of the fort to his son.
  \item \emph{Recall:} say the Hindi for a struggle, then the Hindi for a sacrifice, then say \emph{saumpnā} again
\end{itemize}

\begin{tcolorbox}[breakable,colback=teal!4,colframe=teal!35!black,title={Before you move on}]
What does \hi{सौंपना} mean? (\textquotedblleft{}To entrust; to hand over\textquotedblright{}.) Say it once more.
\end{tcolorbox}
